\uuid{JO54}
\titre{Limite en l'infini}
\auteur{bodin ; hanani ; mzari}
\organisation{exo7}
\licence{CC BY-NC-SA 4.0}
\datecreate{2020-01-20}
\niveau{L1}
\module{Analyse}
\chapitre{Continuité, limite et étude de fonctions réelles}
\sousChapitre{Limite de fonctions}
\difficulte{1}
\theme{Limites des fonctions réelles, limite en l'infini}
\exo7qcmid{lille-1.limites.001}

\contenu{
\begin{question}
\qclassification{123.03}
\qtags{motcle=limite en l'infini}
Soit $f(x)= \frac{x^2+2x+1}{x^2-x-1}$. Quelles sont les assertions vraies ?
\begin{answers}
    \bad{$\lim_{x\to +\infty} f(x)=-1$}

    \good{$\lim_{x\to +\infty} f(x)=1$}

    \good{$\lim_{x\to -\infty} f(x)=1$}

    \bad{$\lim_{x\to -\infty} f(x)=-1$}
\end{answers}

\begin{explanations}
La limite en l'infini d'une fraction rationnelle est la limite de la fraction des monômes de plus haut degré.
\end{explanations}
\end{question}
}
